% Proyecto 2.1: Explorando otras aplicaciones - intérprete SQL.
% Documento autónomo para Overleaf. Compilar informe.tex con pdfLaTeX.
\documentclass[11pt,a4paper]{article}
\usepackage[utf8]{inputenc}
\usepackage[T1]{fontenc}
\usepackage{lmodern}
\usepackage{babel}
\babelprovide[main,import]{spanish}
\usepackage[margin=2.5cm,headheight=18pt,headsep=16pt]{geometry}
\usepackage{graphicx}
\usepackage{xcolor}
\usepackage{booktabs}
\usepackage{longtable}
\usepackage{array}
\usepackage{tabularx}
\usepackage{fancyvrb}
\usepackage{float}
\usepackage[hidelinks,hypertexnames=false]{hyperref}

\definecolor{verdeUVG}{RGB}{0,105,60}
\definecolor{bordeCodigo}{RGB}{190,198,208}
\definecolor{grisTexto}{RGB}{83,91,101}
\newcolumntype{L}[1]{>{\raggedright\arraybackslash}p{#1}}
\newcolumntype{Y}{>{\raggedright\arraybackslash}X}
\DefineVerbatimEnvironment{sqlcode}{Verbatim}{samepage=true,fontsize=\footnotesize,
  frame=single,rulecolor=\color{bordeCodigo},framesep=5pt,
  numbers=left,numbersep=6pt,xleftmargin=12pt,baselinestretch=1.02}
\DefineVerbatimEnvironment{archivo}{Verbatim}{samepage=true,fontsize=\footnotesize,
  frame=single,rulecolor=\color{bordeCodigo},framesep=5pt,
  baselinestretch=1.02}
\setlength{\parindent}{0pt}
\setlength{\parskip}{0.45em}
\setlength{\emergencystretch}{2em}
\setlength{\tabcolsep}{5pt}
\renewcommand{\arraystretch}{1.16}
\setcounter{tocdepth}{2}
\hypersetup{pdftitle={Proyecto 2.1 - Interprete SQL con almacenamiento en archivos},
  pdfauthor={Luis Gonzalez y Joel Jaquez},
  pdfsubject={Construccion de Compiladores - Explorando otras aplicaciones}}

% Encabezado equivalente al de los informes anteriores, sin dependencias extra.
\makeatletter
\def\ps@sqlreport{%
  \def\@oddhead{\vbox{\hbox to\textwidth{\small Proyecto 2.1: Intérprete SQL\hfil ANTLR + C++17}\vskip5pt\hrule height0.4pt}}%
  \let\@evenhead\@oddhead
  \def\@oddfoot{\hfil\thepage\hfil}\let\@evenfoot\@oddfoot}
\makeatother
\pagestyle{sqlreport}
\newcommand{\subhead}[1]{\par\medskip\noindent\textbf{#1}\par\nopagebreak\smallskip}
\newcommand{\captura}[3]{%
  \begin{figure}[H]\centering
  \includegraphics[width=\textwidth,height=0.66\textheight,keepaspectratio]{img/#1}
  \caption{#2}\label{#3}
  \end{figure}}

\begin{document}
\renewcommand{\tablename}{Tabla}
\begin{titlepage}
\begin{center}
\vspace*{0.2cm}
\includegraphics[width=0.24\textwidth]{img/uvg-logo.png}
\vspace{0.8cm}

{\Large\scshape Universidad del Valle de Guatemala}\\[0.30cm]
{\large Facultad de Ingeniería}\\[0.20cm]
{\large Departamento de Ciencias de la Computación}\\[0.20cm]
{\large Construcción de Compiladores}

\vspace{1.0cm}
\rule{\textwidth}{0.8pt}
\vspace{0.45cm}

{\Large\bfseries Proyecto 2.1}\\[0.22cm]
{\large Explorando otras aplicaciones}\\[0.45cm]
{\Huge\bfseries Intérprete SQL}\\[0.30cm]
{\LARGE\bfseries con almacenamiento en archivos}

\vspace{0.45cm}
{\small ANTLR 4.13.2 \ $\bullet$ \ C++17 \ $\bullet$ \ AST \ $\bullet$ \ IDE web}\\[0.25cm]
{\small\href{https://github.com/G1LB3T0/Proyecto1_Compis_Generador_Analizador_Lexico/tree/Compiladores2-Proyecto3}{Repositorio: rama Compiladores2-Proyecto3}}
\vspace{0.45cm}
\rule{\textwidth}{0.8pt}

\vspace{0.9cm}
{\large\bfseries Catedrático}\\[0.25cm]
{\large Ing. Pablo Koch}
\vspace{0.75cm}

{\large\bfseries Integrantes}\\[0.3cm]
{\large Luis Gonzalez --- 23353}\\[0.2cm]
{\large Joel Jaquez --- 23369}

\vfill
{\large Guatemala, 31 de octubre de 2026}\\[0.18cm]
{\small Versión del código documentada: \texttt{037ac10}}
\end{center}
\end{titlepage}

\pagenumbering{roman}
\tableofcontents
\clearpage
\listoffigures
\clearpage
\pagenumbering{arabic}

\section{Introducción}

Este informe presenta la implementación del \textbf{intérprete SQL con
almacenamiento en archivos} correspondiente al Proyecto 2.1, \emph{Explorando
otras aplicaciones}. En el repositorio, esta entrega se identifica mediante
la rama \texttt{Compiladores2-Proyecto3}, por ser la tercera aplicación
desarrollada en la serie de proyectos de compiladores.

El objetivo es aplicar análisis léxico, análisis sintáctico, construcción de
un árbol de sintaxis abstracta (AST), análisis semántico e interpretación a
un lenguaje de consultas. Una sentencia válida se ejecuta sobre esquemas,
registros e índices administrados por el propio programa. El sistema
\textbf{no utiliza SQLite, MySQL, PostgreSQL ni otro DBMS} para almacenar los
datos o ejecutar SQL.

El núcleo está escrito en \textbf{C++17}. ANTLR genera el lexer y el parser
desde \texttt{SQL.g4}; un Visitor construye el AST; el analizador semántico
resuelve tablas, columnas y tipos; y el intérprete produce resultados o
modifica archivos de texto. Flask sirve como adaptador entre el ejecutable
nativo y el IDE web. La interfaz permite inspeccionar tanto el resultado
como las fases que lo producen.

La infraestructura se deriva de los proyectos de Compiscript, pero esta
entrega \textbf{interpreta el AST directamente}: el flujo SQL no genera TAC,
temporales ni registros de activación. Se reutilizan las herramientas de
análisis y presentación, mientras que el lenguaje, el catálogo y la
ejecución de consultas son componentes propios de SQL.

La validación actual comprende \textbf{117 pruebas de integración aprobadas}:
76 casos clasificados como válidos y 41 errores esperados. Además de los
requisitos funcionales, se verifican Unicode, precisión numérica,
reapertura de archivos, acceso directo por índices y recuperación frente a
fallos de escritura.

\section{Objetivos}
\subsection{Objetivo general}

Desarrollar un sistema que interprete y ejecute el subconjunto SQL requerido
por el enunciado sobre un mecanismo de almacenamiento propio, utilizando
las etapas de análisis de un compilador para validar las sentencias antes
de operar sobre archivos.

\subsection{Objetivos específicos}
\begin{itemize}
  \item Definir una gramática SQL y generar el lexer y el parser con ANTLR.
  \item Construir un AST independiente del árbol sintáctico concreto y
        decorarlo con tipos y referencias de columnas.
  \item Validar tablas, columnas, tipos, índices, restricciones y
        subconsultas mediante diagnósticos con ubicación.
  \item Implementar DDL, DML, consultas, JOIN, agregaciones y agrupación.
  \item Persistir esquemas, tipos, registros e índices en archivos de texto
        y mantenerlos consistentes ante modificaciones y fallos.
  \item Usar índices en los predicados compatibles y leer directamente
        las filas candidatas del CSV.
  \item Integrar el pipeline en una interfaz que muestre tokens, árbol,
        AST, errores, resultados y archivos de la base de datos.
  \item Aportar pruebas reproducibles y evidencia visual de funcionamiento.
\end{itemize}

\section{Correspondencia con el enunciado}

El enunciado organiza la entrega en ocho incrementos. La tabla siguiente
relaciona cada uno con los componentes implementados y su forma de
verificación. Los ejemplos ejecutables se encuentran en
\path{backend/examples/sql/}.

\begin{longtable}{@{}L{3.1cm}L{6.2cm}L{5.0cm}@{}}
\toprule
\textbf{Entrega} & \textbf{Implementación} & \textbf{Evidencia}\\
\midrule\endfirsthead
\toprule\textbf{Entrega} & \textbf{Implementación} & \textbf{Evidencia}\\
\midrule\endhead
1. Lexer y parser & \texttt{SQL.g4}, código ANTLR, clasificación de tokens y
listeners propios. & Tokens y árbol; errores \texttt{LEX001}/\texttt{SYN001}.\\
2. AST & \texttt{AstBuilder}, nodos de sentencias y expresiones,
serialización independiente. & Pestaña AST y pruebas de precedencia.\\
3. Semántica & Catálogo, ámbitos de consulta, tipos, restricciones y
diagnósticos \texttt{SEM001}--\texttt{SEM020}. & Rechazo previo y
ubicación de errores.\\
4. CREATE, DROP, INSERT & Creación y eliminación de tablas; filas de
\texttt{VALUES}; catálogo persistente. & Ejemplos 01 y 02; CSV y metadatos.\\
5. SELECT y DML & Proyección, FROM, WHERE, AND/OR, ORDER BY, UPDATE y DELETE.
& Ejemplos 03 y 04; resultados y archivos.\\
6. JOIN y agregaciones & INNER JOIN, COUNT, SUM, AVG, MIN, MAX, GROUP BY y
HAVING. & Ejemplo 05; verificaciones por grupo.\\
7. Índices y subconsultas & CREATE/DROP INDEX; mantenimiento DML; IN/NOT IN,
EXISTS/NOT EXISTS y subconsultas escalares. & Ejemplo 06; planes de acceso y
reapertura.\\
8. Integración & CLI JSON, adaptador Flask, IDE y batería de pruebas.
& Interfaz ejecutable y 117 casos.\\
\bottomrule
\caption{Trazabilidad de los incrementos solicitados.}
\end{longtable}

\section{Arquitectura y reutilización}
\subsection{Flujo de procesamiento}

La función pública \texttt{runScript(source, options)} coordina las fases.
\texttt{RunOptions} selecciona el archivo, la carpeta de almacenamiento y
el modo de ejecución o de análisis. El mismo servicio atiende al CLI, las
pruebas y el IDE.

\begin{center}
\fbox{\begin{minipage}{0.88\textwidth}
\centering\small
\textbf{Texto SQL}\par
$\downarrow$\par
\textbf{SQLLexer} $\rightarrow$ tokens y diagnósticos léxicos\par
$\downarrow$\par
\textbf{SQLParser} $\rightarrow$ árbol sintáctico concreto\par
$\downarrow$\par
\textbf{AstBuilder} $\rightarrow$ AST de sentencias y expresiones\par
$\downarrow$\par
\textbf{SemanticAnalyzer} $\leftrightarrow$ copia del catálogo\par
$\downarrow$\par
\textbf{Executor} $\leftrightarrow$ \textbf{Storage / FileTransaction}\par
$\downarrow$\par
\textbf{Archivos de texto y resultados}\par
$\downarrow$\par
JSON $\rightarrow$ Flask $\rightarrow$ IDE web
\end{minipage}}
\end{center}

ANTLR proporciona los mecanismos para reconocer el lenguaje y recorrer
sus árboles generados; las reglas de significado y ejecución pertenecen
a la aplicación \cite{antlr,antlr-cpp}.

Si existe un error léxico o sintáctico, se omiten la semántica y la
ejecución. Si hay errores semánticos en cualquier sentencia del script,
no se ejecuta ninguna. En modo análisis se construyen y validan las
estructuras, pero no se ejecutan sentencias SQL.

\begin{table}[H]\centering\small
\begin{tabularx}{\textwidth}{@{}Yccc@{}}
\toprule
\textbf{Entrada} & \textbf{AST} & \textbf{Semántica} & \textbf{Ejecución}\\
\midrule
Error léxico o sintáctico & omitido & omitida & omitida\\
Error semántico & disponible & rechazada & omitida\\
Script válido, modo análisis & decorado & correcta & omitida\\
Script válido, modo ejecución & decorado & correcta & habilitada\\
\bottomrule
\end{tabularx}
\caption{Separación y omisión de fases del pipeline SQL.}
\end{table}

Un error que depende de los datos se reporta durante la interpretación.
En ese caso se detiene el script: las sentencias anteriores ya confirmadas
permanecen guardadas y las siguientes se omiten. Esta política se distingue
de los errores semánticos, que se detectan antes de comenzar a ejecutar.

\subsection{Qué se conserva y qué cambia respecto a Compiscript}

\begin{table}[H]\centering\small
\begin{tabularx}{\textwidth}{@{}L{3.2cm}YY@{}}
\toprule
\textbf{Componente} & \textbf{Reutilización} & \textbf{Adaptación SQL}\\
\midrule
ANTLR y CMake & Runtime C++ y generación de código. & Nueva gramática y
nuevas clases generadas.\\
Diagnósticos & Listener que recoge errores y conserva ubicación. & Mensajes
en español y códigos de las fases SQL.\\
Árbol y JSON & Serialización de nodos y escape de cadenas. & AST, catálogo,
resultados, planes y archivos.\\
Servicio y CLI & Un núcleo independiente con salida JSON. & \texttt{--db},
\texttt{--analyze}, \texttt{--catalog} y \texttt{--reset}.\\
Flask e IDE & Editor, archivos, pestañas, visor SVG y reporte de pruebas.
& Resultados SQL y explorador del almacenamiento.\\
Semántica y ejecución & Separación entre reconocimiento y significado.
& Catálogo relacional, consultas, índices y persistencia propios.\\
\bottomrule
\end{tabularx}
\caption{Reutilización de infraestructura y componentes específicos de SQL.}
\end{table}

El módulo anterior permanece en \path{backend/compiscript/}. Su compilación
es opcional con \texttt{BUILD\_COMPISCRIPT=ON}; el proyecto predeterminado
compila el intérprete SQL. La existencia del código anterior en el
repositorio no lo incorpora al pipeline de esta entrega.

\subsection{Organización del núcleo}
\begin{longtable}{@{}L{5.4cm}L{8.9cm}@{}}
\toprule\textbf{Archivo o grupo} & \textbf{Responsabilidad}\\\midrule
\endfirsthead
\toprule\textbf{Archivo o grupo} & \textbf{Responsabilidad}\\\midrule\endhead
\path{grammar/SQL.g4} & Definición del lexer y el parser.\\
\path{generated/SQL*.cpp} & Reconocimiento generado por ANTLR.\\
\path{ast.h}, \path{ast_builder.cpp} & Modelo del AST y Visitor que lo construye.\\
\path{ast_serializer.cpp} & AST decorado en JSON para el visor del IDE.\\
\path{semantic_analyzer.cpp} & Resolución de nombres, tipos y restricciones estáticas.\\
\path{catalog.h/.cpp} & Esquemas, columnas, restricciones e índices.\\
\path{value.h/.cpp} & Valores, comparación, conversión, fechas y Unicode.\\
\path{executor.h/.cpp} & Interpretación de DDL, DML, consultas y expresiones.\\
\path{storage.h/.cpp} & CSV, índices, posiciones, caché y bloqueo de archivos.\\
\path{file_transaction.h/.cpp} & Preparación, respaldos, confirmación y recuperación.\\
\path{service.cpp}, \path{main.cpp} & Orquestación y opciones del CLI.\\
\path{json.cpp}, listeners y serializers & Contrato JSON, diagnósticos, tokens y árbol concreto.\\
\bottomrule
\caption{Componentes de \texttt{backend/sql/}; cabeceras en \texttt{include/sql/} y fuentes en \texttt{src/}.}
\end{longtable}

\section{Reconocimiento del lenguaje y construcción del AST}
\subsection{Gramática SQL}

La gramática combinada \texttt{SQL.g4} admite secuencias de sentencias
separadas por punto y coma. Sus alternativas principales corresponden a
\texttt{CREATE TABLE}, \texttt{DROP TABLE}, \texttt{CREATE INDEX},
\texttt{DROP INDEX}, \texttt{INSERT}, \texttt{UPDATE}, \texttt{DELETE} y
\texttt{SELECT}. Se permiten punto y coma final opcional, sentencias vacías
y comentarios de línea o de bloque.

La opción \texttt{caseInsensitive=true} hace que palabras reservadas e
identificadores se reconozcan sin distinguir mayúsculas. El AST normaliza
los nombres a minúsculas; el contenido de las cadenas conserva su texto.
Los tokens reservados tienen prefijo \texttt{K\_}, por ejemplo
\texttt{K\_SELECT} y \texttt{K\_NULL}, para evitar conflictos con macros de
C/C++ y de Windows.

Las reglas distinguen enteros, decimales, cadenas entre comillas simples,
booleanos, NULL y literales DATE. En una cadena, la comilla simple se
escapa duplicándola: \texttt{'O''Brien'}. Los comentarios \texttt{--} y
\texttt{/* ... */} se descartan durante el análisis léxico.

\subsection{Precedencia y árbol sintáctico concreto}

Las alternativas de expresión establecen la precedencia de signos unarios,
multiplicación y división, suma y resta, comparaciones, pruebas de NULL,
pertenencia, NOT, AND y OR. Los paréntesis cambian la agrupación.
Así, \texttt{2 + 3 * 4} produce 14 y \texttt{TRUE OR FALSE AND FALSE}
produce TRUE.

El parser construye el \emph{parse tree}, que conserva las reglas
intermedias y los tokens de puntuación. \texttt{tree\_serializer} lo
convierte en nodos con nombre, texto, categoría, línea, columna e hijos.
Este árbol permite explicar cómo se reconoció una sentencia, pero no es
la estructura que interpreta directamente el ejecutor.

\subsection{AST y decoración semántica}

\texttt{AstBuilder} deriva de \texttt{SQLBaseVisitor}. Cada alternativa
crea nodos con significado y descarta comas, punto y coma y paréntesis de
agrupación. Las sentencias se almacenan en un \texttt{std::variant}; las
expresiones se distinguen por \texttt{ExprKind}.

\begin{table}[H]\centering\small
\begin{tabularx}{\textwidth}{@{}L{3.6cm}Y@{}}
\toprule\textbf{Tipo de expresión} & \textbf{Ejemplos}\\\midrule
Literal y columna & \texttt{18}, \texttt{'Ana'}, \texttt{TRUE}, \texttt{e.nombre}.\\
Unaria y binaria & \texttt{-edad}, \texttt{NOT activo}, \texttt{edad >= 18}.\\
Pruebas y pertenencia & \texttt{IS NULL}, \texttt{IN (...)}, \texttt{NOT IN (SELECT ...)}.\\
Subconsulta & EXISTS y subconsulta escalar.\\
Agregación & COUNT, SUM, AVG, MIN y MAX, con DISTINCT cuando corresponde.\\
\bottomrule
\end{tabularx}
\caption{Familias de expresiones representadas en el AST.}
\end{table}

Para \texttt{SELECT nombre FROM estudiantes WHERE edad >= 18}, el AST
contiene una proyección, una tabla y un predicado de comparación. La
semántica agrega el tipo de cada expresión y, para las columnas, una
posición \texttt{slot} en la fila combinada y una profundidad
\texttt{depth}. El ejecutor consulta estas posiciones en lugar de volver
a resolver cada nombre por texto.

Los nodos conservan la ubicación y el fragmento fuente. Por ello el AST
puede servir tanto para ejecutar como para presentar diagnósticos y
explicar las decisiones en el IDE.

\section{Análisis semántico y manejo de errores}
\subsection{Catálogo y ámbitos de consulta}

El catálogo sustituye a la tabla de símbolos de un lenguaje de programación
general: describe tablas, columnas, tipos, restricciones e índices.
\texttt{SemanticAnalyzer} trabaja con una copia del catálogo y procesa
el script en orden. Un CREATE válido incorpora su esquema a esa copia,
de modo que las sentencias siguientes pueden referirse a la tabla sin
que todavía se haya escrito ningún archivo.

Cada SELECT abre un ámbito con las columnas de FROM y de los JOIN ya
incorporados. El predicado ON solo puede usar las tablas disponibles en
ese punto. Un nombre sin calificador que coincide con varias columnas
produce una referencia ambigua. Los alias distinguen las tablas y sus usos.

En una subconsulta se abre otro ámbito. La búsqueda comienza en la
consulta interior y continúa hacia los niveles externos. Si una columna
se resuelve en un nivel exterior, se registra \texttt{depth > 0} y la
subconsulta se marca como correlacionada.

\subsection{Tipos, conversiones y restricciones}

\begin{table}[H]\centering\small
\begin{tabularx}{\textwidth}{@{}L{2.8cm}Y@{}}
\toprule\textbf{Tipo} & \textbf{Representación y validación}\\\midrule
INT & Entero con signo de 64 bits; conversiones y operaciones comprobadas frente a desbordamientos.\\
FLOAT & Punto flotante \texttt{double}; admite promoción desde INT y persistencia decimal que reconstruye el mismo valor binario.\\
VARCHAR(n) & Texto UTF-8 con máximo de \(n\) puntos de código Unicode. Se preservan comillas, comas y saltos de línea.\\
BOOLEAN & TRUE o FALSE; los operadores lógicos exigen operandos compatibles.\\
DATE & Formato YYYY-MM-DD; comprueba mes, día y años bisiestos. Un literal de texto válido puede convertirse al asignar o comparar con DATE.\\
\bottomrule
\end{tabularx}
\caption{Tipos mínimos solicitados y decisiones de implementación.}
\end{table}

NULL representa ausencia de valor y no constituye un sexto tipo de columna.
Las comparaciones ordinarias con NULL producen un resultado desconocido;
WHERE solo conserva filas cuyo predicado resulta TRUE. Se implementan
AND y OR con lógica de tres valores y cortocircuito.

Las restricciones adicionales incluyen PRIMARY KEY, NOT NULL, UNIQUE y
FOREIGN KEY. Las referencias deben apuntar a columnas PRIMARY KEY o
UNIQUE de tipo compatible. Se emplea la política RESTRICT para impedir
borrar o cambiar una clave que otra tabla referencia.

Las restricciones que dependen de valores almacenados se comprueban en
ejecución. Por ejemplo, una columna puede ser válida en el esquema y,
sin embargo, un INSERT fallar porque repite una clave o referencia un
registro inexistente.

\subsection{Códigos y ubicación de los diagnósticos}

Los diagnósticos contienen código, categoría, severidad, mensaje, número de
sentencia, línea y columna. Las posiciones se muestran al usuario desde 1.
El IDE permite seleccionar un error para ir al punto correspondiente del
editor. Los listeners propios sustituyen a los mensajes de consola de
ANTLR y producen \texttt{LEX001} o \texttt{SYN001}.

\begin{table}[H]\centering\small
\begin{tabularx}{\textwidth}{@{}L{2.6cm}Y@{}}
\toprule\textbf{Código} & \textbf{Validación semántica}\\\midrule
SEM001--004 & Tabla inexistente; columna inexistente; referencia ambigua;
tabla o alias inválido.\\
SEM005--009 & Tipos incompatibles; cantidad incorrecta de valores; NULL u
omisión en columna obligatoria; exceso de VARCHAR; fecha inválida.\\
SEM010--012 & Tabla ya existente; columna duplicada; restricción inválida.\\
SEM013--015 & Índice inexistente; índice duplicado; subconsulta IN o escalar
con una cantidad de columnas distinta de una.\\
SEM016--017 & Uso inválido de agregación o columna no agrupada; intento de
eliminar una tabla referenciada.\\
SEM018--020 & Literal fuera de rango; condición no booleana; posición de
ORDER BY fuera del resultado.\\
\bottomrule
\end{tabularx}
\caption{Agrupación de los veinte códigos semánticos implementados.}
\end{table}

Los errores de ejecución son \texttt{EXE001} (PRIMARY KEY), \texttt{EXE002}
(UNIQUE), \texttt{EXE003} (FK sin padre), \texttt{EXE004} (registro
referenciado), \texttt{EXE005} (subconsulta escalar con varias filas),
\texttt{EXE006} (división entre cero), \texttt{EXE007} (valor o resultado
fuera del tipo admitido) y \texttt{EXE008} (almacenamiento o bloqueo).

Una subconsulta escalar con más de una columna se rechaza durante la
semántica. Su cantidad de filas se comprueba al evaluarla, porque depende
del contenido de la base y de sus predicados. Si retorna cero filas,
su resultado escalar es NULL.

\section{Diseño del almacenamiento en archivos}
\subsection{Distribución de los archivos}

\begin{archivo}
database/
  metadata.txt
  estudiantes.csv
  estudiantes.rows
  cursos.csv
  cursos.rows
  indexes/
    pk_estudiantes.idx
    idx_estudiantes_edad.idx
  .lock
  .transaction/           (solo durante una confirmacion)
\end{archivo}

Todos los datos de la base son archivos de texto. \texttt{metadata.txt}
describe los esquemas; cada tabla tiene un CSV; cada índice, un archivo
\texttt{.idx}; y \texttt{.rows} relaciona los identificadores de fila con
posiciones físicas en el CSV. El bloqueo y el diario son mecanismos de
coordinación del almacenamiento, no un servicio de base de datos externo.

\subsection{Catálogo: esquemas, tipos y restricciones}

El catálogo utiliza instrucciones de texto por línea. Para cada tabla
conserva el siguiente rowid libre y sus columnas con tipos y restricciones.
Los índices se describen con su tabla, columna y propiedades.

\begin{archivo}
TABLE estudiantes
  NEXT_ROWID 9
  COLUMN id INT PRIMARY_KEY NOT_NULL
  COLUMN nombre VARCHAR(50) NOT_NULL
  COLUMN edad INT
  COLUMN promedio FLOAT
  COLUMN activo BOOLEAN
  COLUMN ingreso DATE
  COLUMN carrera_id INT REFERENCES carreras.id
END

INDEX pk_estudiantes ON estudiantes.id UNIQUE AUTO
INDEX idx_estudiantes_edad ON estudiantes.edad
\end{archivo}

La marca AUTO identifica los índices creados por PRIMARY KEY o UNIQUE.
No se permite eliminarlos mediante DROP INDEX, porque sostienen la
restricción que los originó. El catálogo se lee al abrir la base y se
serializa dentro de la misma confirmación que sus archivos afectados.

\subsection{Registros CSV y valores especiales}

Cada registro comienza con \texttt{\_\_rowid}, un identificador interno
estable que no cambia cuando se eliminan otras filas. Los índices apuntan
a este identificador y no al ordinal de la fila dentro del archivo.

\begin{archivo}
__rowid,id,nombre,edad,promedio,activo,ingreso,carrera_id
1,1,"Ana López",20,91.5,TRUE,2023-01-15,1
4,4,"Pedro Méndez",18,,TRUE,2024-01-10,
\end{archivo}

El fragmento anterior ilustra el formato. Todo texto se escribe entre
comillas dobles y las comillas
internas se duplican. Un campo vacío sin comillas representa NULL y
\texttt{""} representa una cadena vacía. El lector respeta comas y saltos
de línea dentro de campos entrecomillados.

FLOAT se escribe mediante una representación decimal suficientemente
precisa para reconstruir el mismo \texttt{double}. La conversión de
entrada controla el rango y rechaza valores no finitos. VARCHAR cuenta
puntos de código de UTF-8: por ejemplo, \texttt{José} ocupa cuatro
caracteres y puede almacenarse en VARCHAR(4).

\subsection{Posiciones de fila y lectura directa}

Un índice entrega rowids. Para evitar recorrer el CSV hasta encontrarlos,
el archivo \texttt{<tabla>.rows} conserva la posición y longitud de cada
registro en bytes:

\begin{archivo}
ROWS 1 <tamano_del_csv_en_bytes>
<rowid> <offset_en_bytes> <longitud_en_bytes>
\end{archivo}

\texttt{readRows} abre el archivo en modo binario, hace \texttt{seekg} a
cada posición candidata y decodifica únicamente ese registro. Las
posiciones incluyen la representación CSV completa, sus comillas y los
saltos de línea internos; por eso se calculan al serializar la tabla.
INSERT, UPDATE y DELETE confirman el CSV y sus nuevas posiciones juntos.

Al abrir un directorio de posiciones se comprueban su versión, el tamaño
del CSV y los límites de sus entradas. Una base antigua sin \texttt{.rows}
sigue funcionando: se carga la tabla y se calculan las posiciones en
memoria; la siguiente modificación las guarda en disco.

\subsection{Confirmación, recuperación y concurrencia}

Las modificaciones trabajan sobre copias de la tabla y sus índices.
Superadas las restricciones, una sentencia prepara un conjunto de
escrituras y eliminaciones. \texttt{FileTransaction} realiza estos pasos:

\begin{enumerate}
  \item Registra en un manifiesto todos los archivos afectados y si ya existían.
  \item Respalda los originales y termina de escribir todos los temporales.
        Se comprueba también el cierre de cada archivo.
  \item Marca el diario como preparado y comienza a instalar los archivos.
  \item Después de instalar todos los cambios, coloca el marcador de
        confirmación mediante un renombrado.
  \item Elimina los temporales y respaldos que ya no se necesitan.
\end{enumerate}

Si falla la preparación, ningún original se modifica. Si falla o se
interrumpe la instalación, el diario permite restaurar los respaldos y
eliminar archivos nuevos no confirmados. Al volver a abrir la base se
recupera una operación pendiente antes de cargar el catálogo. Un diario
marcado como confirmado solo necesita limpieza; no se revierte.

Un bloqueo del sistema operativo protege la ejecución completa del CLI.
Mientras otro proceso conserva el bloqueo, una nueva ejecución recibe
\texttt{EXE008}; así no puede recuperar un diario que todavía está activo.
El bloqueo se libera al cerrar el proceso, incluso tras una terminación
inesperada.

La unidad de confirmación es una sentencia, no el script completo. No se
implementan transacciones SQL de varias sentencias ni durabilidad frente
a cortes eléctricos mediante \texttt{fsync}. La recuperación descrita
cubre errores de escritura e interrupciones del proceso.

\section{Interpretación y ejecución de consultas}
\subsection{DDL, INSERT, UPDATE y DELETE}

CREATE TABLE genera el esquema, el CSV vacío, las posiciones y los índices
automáticos. DROP TABLE elimina la entrada del catálogo y los archivos
asociados, siempre que ninguna otra tabla la referencie. CREATE INDEX
construye un índice a partir de los datos existentes; si es UNIQUE,
rechaza claves repetidas antes de confirmarlo.

INSERT evalúa cada fila de VALUES, convierte sus valores y comprueba
NOT NULL, unicidad y claves foráneas. Las columnas omitidas reciben NULL
cuando el esquema lo permite. Si cualquier fila falla, la sentencia
entera se descarta.

UPDATE localiza sus candidatos, evalúa todas las asignaciones sobre la
fila original y valida los resultados. Retira primero las claves
anteriores de los índices y luego incorpora las nuevas. Esto permite,
por ejemplo, desplazar varias claves sin detectar duplicados transitorios
que no existirán en el estado final.

DELETE evalúa WHERE y verifica las referencias externas antes de retirar
filas y entradas de índices. Los cambios se confirman junto con las
posiciones físicas correspondientes.

\subsection{Orden de evaluación de SELECT}

El orden lógico implementado es el siguiente:

\begin{archivo}
FROM -> INNER JOIN -> WHERE
     -> GROUP BY / agregaciones -> HAVING
     -> proyeccion SELECT -> DISTINCT -> ORDER BY -> LIMIT
\end{archivo}

WHERE se evalúa sobre cada fila candidata, incluso cuando se utilizó un
índice para reducir el conjunto. La proyección produce las expresiones
solicitadas y expande \texttt{*} o \texttt{tabla.*} a columnas concretas.
ORDER BY admite expresiones, alias de salida y posiciones válidas.

\subsection{JOIN, agrupación y agregaciones}

INNER JOIN combina las filas que cumplen ON. Si una columna indexada de
la tabla derecha se compara con un valor ya conocido de la izquierda,
se emplea un \emph{Index Nested Loop}; en otro caso se prueban las
combinaciones mediante \emph{Nested Loop}. Cada unión amplía la fila
combinada y conserva las posiciones resueltas durante la semántica.

GROUP BY construye grupos por el valor de sus expresiones. COUNT(*) cuenta
filas; COUNT(expr) cuenta valores no nulos. SUM y AVG exigen argumentos
numéricos; MIN y MAX seleccionan extremos de tipos admitidos. Las
agregaciones ignoran NULL en sus argumentos. Sin filas, COUNT produce 0
y las demás agregaciones producen NULL.

En una consulta agregada, las columnas fuera de funciones deben estar
agrupadas. HAVING filtra grupos después de las agregaciones. Su presencia
sin GROUP BY crea un grupo implícito, incluso sin funciones agregadas:
\texttt{SELECT 1 HAVING FALSE;} devuelve cero filas.

\subsection{Subconsultas y lógica de NULL}

IN y NOT IN comparan una expresión con una lista o con la columna de una
subconsulta. EXISTS y NOT EXISTS verifican si esta produce filas. Una
subconsulta escalar debe tener una columna y, al evaluarse, como máximo
una fila. Estos conceptos se contrastaron con documentación SQL pública
como referencia semántica; esa consulta documental no incorpora ningún
DBMS al programa \cite{subconsultas}.

Las subconsultas no correlacionadas se calculan una vez por sentencia y
se guardan en caché. Las correlacionadas usan las columnas externas y se
reevalúan para cada contexto externo necesario. El plan registra la
correlación, la cantidad de ejecuciones y el plan de la primera evaluación.

Si IN no encuentra coincidencia y existe NULL en el conjunto, el resultado
es desconocido. El conjunto vacío tiene un tratamiento distinto:
\texttt{NULL IN (subconsulta vacía)} es FALSE y su negación es TRUE.
Esta distinción evita perder filas con \texttt{NOT IN} cuando la
subconsulta no produce resultados.

\section{Índices y planificación del acceso}

Cada índice se representa en memoria mediante un
\texttt{std::map<Value, std::set<rowid>>}. Es una estructura ordenada con
búsquedas logarítmicas sobre sus claves; no es un árbol B almacenado por
páginas. El archivo \texttt{.idx} es texto ordenado y se carga cuando
se necesita.

\begin{archivo}
# INDEX idx_estudiantes_edad ON estudiantes(edad)
# valor,rowids (ordenado por valor)
16,8
17,2
18,4
19,5
20,1
21,7
22,3
25,6
\end{archivo}

Las claves NULL se omiten. Los comentarios pertenecen únicamente al
encabezado: una línea que comienza con \texttt{\#} dentro de una cadena
multilínea forma parte de la clave y debe conservarse.

El planificador reconoce igualdades, IN no negado y comparaciones de rango,
además de \texttt{<>}, cuando el otro operando puede evaluarse antes de
leer la tabla objetivo. También aprovecha valores de tablas anteriores
en un JOIN o de consultas externas.

\begin{table}[H]\centering\small
\begin{tabularx}{\textwidth}{@{}L{4.0cm}Y@{}}
\toprule\textbf{Predicado} & \textbf{Acceso}\\\midrule
\texttt{edad >= 18} & Busca las claves del rango y obtiene sus rowids.\\
\texttt{id IN (1,3)} & Une los rowids encontrados para cada valor.\\
\texttt{id = 1 OR id = 3} & Une ambos accesos si las dos ramas pueden usar índices.\\
\texttt{id >= 2 AND edad >= 19} & Intersecta los accesos disponibles; evalúa el predicado completo al final.\\
OR con una rama no indexable & Recorre la tabla para no excluir resultados de esa rama.\\
\bottomrule
\end{tabularx}
\caption{Decisiones del planificador para predicados indexables.}
\end{table}

INSERT agrega claves, UPDATE retira las anteriores e incorpora las nuevas
y DELETE elimina sus referencias. El archivo de posiciones permite
traducir los rowids a lecturas directas del CSV. El mismo mecanismo se
utiliza en los JOIN indexados.

\subhead{Evidencia de lectura selectiva}
La prueba de una consulta fría por PRIMARY KEY sobre estudiantes verifica
que el almacenamiento decodifica una fila. Otra prueba combina tres
tablas con filtros e índices y verifica tres filas leídas en total.
No se comprueba únicamente la presencia de una etiqueta en el plan.

El contador \texttt{database.reads.table\_rows} registra filas decodificadas
por el almacenamiento. No equivale a tiempo de ejecución ni al total de
bytes del sistema operativo: también se cargan las estructuras de índices
y posiciones, y el explorador del IDE lee una vista previa acotada de
los archivos. Las modificaciones continúan reescribiendo la tabla completa.

\clearpage
\section{Interfaz gráfica y evidencia de funcionamiento}

El servidor Flask ejecuta \texttt{sql\_cli} con la carpeta de base de datos
y el SQL recibido; devuelve al navegador el JSON del núcleo. Las rutas
\texttt{/api/run}, \texttt{/api/database} y \texttt{/api/tests} conectan
ejecución, inspección y pruebas. La separación entre rutas, plantillas y
archivos estáticos sigue la organización habitual de Flask \cite{flask}.

Las capturas siguientes proceden de la aplicación real, ejecutada con
una base de demostración aislada. Se cargaron los ejemplos 01 y 02 y se
crearon los índices de edad y estudiante de inscripción. Los datos
contienen 3 carreras, 8 estudiantes, 6 cursos y 13 inscripciones.
Las imágenes muestran detalles de los paneles, con zoom para facilitar
su lectura; no son maquetas.

\subsection{Consulta filtrada y plan de ejecución}

La consulta del enunciado se amplió para mostrar edad y ordenar el
resultado. Devuelve seis estudiantes: Pedro (18), María (19), Ana (20),
Lucía (21), Sofía (22) y Jorge (25). El plan utiliza
\texttt{idx\_estudiantes\_edad} y el contador registra seis filas
decodificadas de las ocho existentes.

\captura{general.png}{SELECT con WHERE, ORDER BY, resultados y acceso por índice de edad.}{fig:consulta}

\clearpage
\subsection{Tokens}

La tabla de tokens expone categoría, tipo ANTLR, texto y ubicación. Permite
verificar que una palabra reservada, un identificador o un operador
llegaron correctamente al parser.

\captura{tokens.png}{Tokens de la consulta SELECT nombre FROM estudiantes WHERE edad >= 18.}{fig:tokens}

\clearpage
\subsection{Árbol sintáctico concreto}

El árbol sintáctico conserva las reglas de la gramática y sus hojas.
El visor permite expandir, contraer y cambiar el zoom para recorrer
la estructura reconocida por ANTLR.

\captura{arbol.png}{Árbol sintáctico concreto de una consulta, en modo de análisis.}{fig:arbol}

\clearpage
\subsection{AST visible y tipos resueltos}

El AST representa la misma entrada mediante SELECT, proyección, FROM y
WHERE. Las columnas ya muestran sus tipos. Esta vista permite relacionar
la consulta escrita con la estructura utilizada por la semántica y el
intérprete, sin los niveles auxiliares del árbol concreto.

\captura{ast.png}{AST decorado con la proyección, la tabla y el predicado de comparación.}{fig:ast}

\clearpage
\subsection{JOIN, agregaciones y HAVING}

La consulta une cursos e inscripciones, cuenta las inscripciones por
curso y calcula la nota máxima. HAVING conserva los cursos con al menos
dos inscripciones; ORDER BY utiliza el alias del conteo. Se obtienen
cinco grupos. Cálculo tiene cuatro inscripciones y Compiladores tres.

\captura{join.png}{INNER JOIN con COUNT, MAX, GROUP BY, HAVING y ORDER BY.}{fig:join}

\clearpage
\subsection{Subconsultas con IN}

El ejemplo con IN retorna a Ana López, Luis Pérez, María García y Sofía
Castillo: sus identificadores aparecen entre las inscripciones del
curso 3. La subconsulta se calcula una vez y sus valores se buscan en
el índice de la PRIMARY KEY de estudiantes.

\begin{sqlcode}
SELECT nombre
FROM estudiantes
WHERE id IN (SELECT estudiante_id
             FROM inscripciones WHERE curso_id = 3)
ORDER BY nombre;
\end{sqlcode}

\captura{in.png}{IN con SELECT anidado y acceso al índice de estudiantes.}{fig:in}

\clearpage
\subsection{EXISTS correlacionado y subconsulta escalar}

En el caso correlacionado se buscan estudiantes con alguna nota de al
menos 95. Solo Lucía Herrera cumple. La subconsulta usa \texttt{e.id}
del contexto externo y el plan registra ocho evaluaciones, una por
estudiante, con búsquedas por el índice de inscripciones.

\captura{exists.png}{EXISTS correlacionado con una referencia a la fila externa.}{fig:exists}

La variante escalar compara el promedio individual con el promedio global
de 82.75. Devuelve Lucía (95), Ana (91.5), María (88) y Sofía (85.25).
El código y el resultado reproducible de esta variante se incluyen en
\path{evidencia/sql/escalar.sql} y
\path{evidencia/consultas-resultados.json}.

\clearpage
\subsection{Catálogo e índices persistidos}

La pestaña Base de datos muestra las columnas, restricciones, cantidad
de filas e índices de cada tabla, junto con el contenido de los archivos.
El archivo elegido en la figura siguiente es el índice de edad: sus
claves y rowids corresponden al ejemplo descrito en el diseño.

\captura{indice.png}{Explorador del almacenamiento con el contenido real de un índice .idx.}{fig:indice}

\clearpage
\subsection{Rechazo de entradas inválidas}

La primera captura reúne tres errores estáticos independientes: una tabla
inexistente, una columna inexistente y un valor de INSERT incompatible.
El script no se ejecuta. La segunda muestra un error que depende de la
evaluación: división entre cero, reportada como \texttt{EXE006}.

\captura{semantica.png}{Errores semánticos con código, sentencia y ubicación en el editor.}{fig:semantica}
\captura{ejecucion.png}{Error de ejecución por división entre cero.}{fig:ejecucion}

\clearpage
\section{Pruebas y resultados de validación}
\subsection{Metodología}

\texttt{sql\_tests} invoca el mismo servicio C++ que utiliza el IDE. Cada
caso tiene una carpeta temporal aislada; cuando corresponde, se carga un
esquema de ejemplo. La verificación compara filas, tipos, diagnósticos,
planes, contadores y contenido de archivos. Algunos casos realizan una
segunda ejecución para comprobar persistencia.

Un error esperado aprueba cuando el intérprete rechaza la entrada con
el código requerido. El total de 76 casos válidos corresponde a la
clasificación del reporte: incluye también verificaciones de omisión de
fases y de recuperación que examinan respuestas fallidas deliberadas.
Los 41 errores esperados no son fallos del sistema.

\begin{table}[H]\centering\small
\begin{tabularx}{\textwidth}{@{}Yrr@{}}
\toprule\textbf{Categoría} & \textbf{Casos} & \textbf{Aprobados}\\\midrule
Lexer y parser & 7 & 7\\
AST & 4 & 4\\
Análisis semántico & 26 & 26\\
CREATE, DROP e INSERT & 10 & 10\\
SELECT, WHERE, UPDATE y DELETE & 12 & 12\\
JOIN, agregaciones y GROUP BY & 7 & 7\\
Índices & 10 & 10\\
Subconsultas & 7 & 7\\
Almacenamiento en archivos & 3 & 3\\
Regresiones de robustez & 31 & 31\\\midrule
\textbf{Total} & \textbf{117} & \textbf{117}\\
\bottomrule
\end{tabularx}
\caption{Resultados de la batería ejecutada para esta versión del informe.}
\end{table}

\subsection{Regresiones incorporadas}

Las 31 regresiones adicionales se organizan en dos conjuntos: 16 casos
de tipos, expresiones y texto indexado; y 15 de planificación,
almacenamiento, recuperación y concurrencia. Entre las verificaciones
principales se encuentran las siguientes:

\begin{itemize}
  \item VARCHAR con acentos y caracteres de varios bytes, tanto en
        validación semántica como en valores calculados por UPDATE.
  \item Literales FLOAT fuera de rango, valores subnormales, precisión
        después de reabrir y desbordamientos de aritmética y SUM.
  \item HAVING con grupo implícito y NULL frente a subconsultas vacías.
  \item Claves indexadas con saltos de línea, líneas vacías y \texttt{\#}.
  \item Unión e intersección de índices y recorrido correcto de un OR con
        una rama no indexable.
  \item Filas efectivamente decodificadas y actualización de posiciones
        tras INSERT, UPDATE y DELETE.
  \item Errores de escritura durante INSERT, UPDATE, CREATE y DROP sin
        confirmar cambios parciales.
  \item Recuperación de una instalación interrumpida y conservación de
        una confirmación terminada cuyo diario requiere limpieza.
  \item Rechazo de una segunda ejecución sobre una base bloqueada.
\end{itemize}

También se ejecutó la batería con AddressSanitizer y
UndefinedBehaviorSanitizer: los 117 casos aprobaron sin diagnósticos de
esas herramientas. Esta comprobación complementa las pruebas funcionales;
no constituye una demostración de ausencia de errores en entradas no
cubiertas.

\captura{pruebas.png}{Reporte de la batería del intérprete SQL: 117 aprobadas y cero fallos reales.}{fig:pruebas}

\subsection{Evidencia reproducible}

La carpeta del informe conserva \path{evidencia/pruebas.json}, con el
reporte íntegro, y \path{evidencia/consultas-resultados.json}, con SQL,
resultados, diagnósticos, planes y contadores de las demostraciones.
\path{evidencia/sql/00_preparar_base.sql} reproduce el esquema, los datos
y los índices usados en las capturas. Cada consulta dispone de un archivo
SQL independiente para que las demostraciones de error no impidan ejecutar
las consultas válidas.

\section{Trazabilidad del desarrollo}

La rama del proyecto contiene seis commits de implementación inicial y
seis de refuerzo. En conjunto, los doce commits descritos tienen una
distribución de seis por integrante. Las identidades coinciden con las
utilizadas en las entregas anteriores.

\begin{table}[H]\centering\small
\begin{tabularx}{\textwidth}{@{}L{2.5cm}L{5.4cm}Y@{}}
\toprule\textbf{Usuario} & \textbf{Correo} & \textbf{Integrante}\\\midrule
\texttt{G1LB3T0} & \texttt{luis.yo09@hotmail.com} & Luis Gonzalez, 23353\\
\texttt{jaq23369} & \texttt{jaq23369@uvg.edu.gt} & Joel Jaquez, 23369\\
\bottomrule
\end{tabularx}
\caption{Identidades registradas en Git.}
\end{table}

\begin{longtable}{@{}L{1.8cm}L{2.3cm}L{10.2cm}@{}}
\toprule\textbf{Commit} & \textbf{Autor} & \textbf{Aporte}\\\midrule\endfirsthead
\toprule\textbf{Commit} & \textbf{Autor} & \textbf{Aporte}\\\midrule\endhead
\texttt{68ae80c} & G1LB3T0 & Gramática SQL y generación del parser ANTLR.\\
\texttt{0426696} & jaq23369 & Construcción del AST y análisis semántico SQL.\\
\texttt{2016eb1} & G1LB3T0 & Almacenamiento en archivos, intérprete e integración del CLI.\\
\texttt{f0ad7c3} & G1LB3T0 & Batería inicial de 86 pruebas de integración.\\
\texttt{ef55839} & jaq23369 & Adaptación del IDE, endpoints y ejemplos SQL.\\
\texttt{510c384} & jaq23369 & Guía de arquitectura, reutilización y defensa.\\
\texttt{2f7be06} & G1LB3T0 & Unicode, tipos numéricos, HAVING y condiciones con NULL.\\
\texttt{bdd28f8} & jaq23369 & Confirmación conjunta, recuperación y bloqueo de archivos.\\
\texttt{10056ad} & G1LB3T0 & Acceso por posiciones y combinación de índices.\\
\texttt{23b9d6b} & jaq23369 & Dieciséis regresiones de tipos y expresiones.\\
\texttt{c0a4b01} & G1LB3T0 & Quince regresiones de acceso, disco y recuperación; 117 casos.\\
\texttt{037ac10} & jaq23369 & Documentación del almacenamiento y de su validación.\\
\bottomrule
\caption{Historial funcional del proyecto SQL, en orden cronológico.}
\end{longtable}

\section{Compilación y ejecución}

Se requiere un compilador C++17, CMake 3.20 o superior, Java para generar
ANTLR y Python con Flask para el IDE. Los archivos C++ generados ya están
en el repositorio; los scripts de construcción también permiten
regenerarlos desde la gramática.

\begin{archivo}
# macOS o Linux
./scripts/build_sql.sh

# Windows (PowerShell)
.\scripts\build_sql.ps1

# Pruebas de integracion
ctest --test-dir build/sql --output-on-failure

# Ejecucion desde consola
./sql_cli --db database --file ejemplo.sql
./sql_cli --db database --analyze --file ejemplo.sql
./sql_cli --db database --catalog

# IDE
python3 app.py
# Abrir http://127.0.0.1:5050
\end{archivo}

Si el runtime de ANTLR ya se encuentra descargado, la variable
\texttt{ANTLR4\_LOCAL\_SOURCE} permite construir sin volver a descargarlo.
El CLI devuelve un documento JSON y utiliza código de salida 0 para
éxito, 1 para errores del pipeline y 2 para problemas de uso o lectura
de la entrada.

Para reproducir las figuras en otra base de datos se recomienda utilizar
una carpeta nueva y ejecutar primero el script de preparación de
\texttt{evidencia/sql/}. Los scripts de error deben ejecutarse por
separado, puesto que un error semántico impide ejecutar todo su script.

\section{Alcance, límites y mejoras futuras}

El proyecto implementa un subconjunto SQL orientado a los requisitos del
curso. Sus límites deben diferenciarse de los comportamientos comprobados:

\begin{itemize}
  \item Se implementa INNER JOIN. No se incluyen LEFT JOIN, consultas
        UNION, vistas ni procedimientos almacenados.
  \item Los índices, las claves primarias y las claves foráneas son de
        una columna. No hay un índice B+ por páginas ni un optimizador
        basado en estimaciones de costo.
  \item Los índices y las posiciones se cargan en memoria. Las lecturas
        indexadas del CSV son selectivas; los recorridos completos y las
        modificaciones cargan la tabla y pueden reescribirla completa.
  \item La colación de texto es binaria. VARCHAR mide puntos de código,
        sin normalización Unicode ni conteo de grafemas compuestos.
  \item Las claves foráneas utilizan RESTRICT; no se implementa CASCADE.
  \item La confirmación cubre una sentencia. No hay BEGIN/COMMIT SQL de
        varias sentencias ni garantía frente a cortes eléctricos con fsync.
  \item El bloqueo exclusivo evita ejecuciones simultáneas sobre la misma
        carpeta; no ofrece concurrencia de múltiples lectores y escritores.
\end{itemize}

Como evolución, sería posible incorporar índices compuestos y por páginas,
actualización localizada de registros, planificación por costo y un
protocolo de durabilidad con sincronización de archivos. Estas mejoras
no forman parte de los resultados afirmados por esta entrega.

\section{Conclusiones}

\begin{enumerate}
  \item Las técnicas de construcción de compiladores permiten interpretar
        un lenguaje de consultas completo dentro del alcance del curso:
        reconocimiento, AST, resolución semántica e interpretación quedan
        separados y son observables en la interfaz.
  \item El almacenamiento propio representa esquemas, tipos, restricciones,
        filas e índices mediante archivos de texto. El catálogo permite
        validar el SQL y el diario mantiene la coherencia de las
        modificaciones ante los fallos cubiertos.
  \item Los índices se utilizan para obtener candidatos y las posiciones
        persistidas permiten leerlos directamente del CSV. La decisión de
        acceso se comprueba con resultados y contadores de filas, además
        del plan mostrado al usuario.
  \item La reutilización de ANTLR, el contrato JSON, Flask y el visor del
        IDE permitió concentrar el trabajo nuevo en el lenguaje SQL y su
        ejecución. El resultado cumple un flujo de interpretación distinto
        de la generación de código intermedio de Compiscript.
  \item Las 117 pruebas y las demostraciones reproducibles aportan evidencia
        del comportamiento funcional y de los casos límite examinados.
        La interfaz y el informe permiten seguir una sentencia desde sus
        tokens hasta los resultados o los archivos modificados.
\end{enumerate}

\begin{thebibliography}{9}
\raggedright
\bibitem{enunciado}
Construcción de Compiladores. \emph{Proyecto 2.1: Intérprete SQL con
Almacenamiento en Archivos}. Enunciado del curso, septiembre de 2026.

\bibitem{antlr}
ANTLR. \emph{ANTLR: herramienta de reconocimiento de lenguajes}.
\url{https://www.antlr.org/}. Consulta: 31 de octubre de 2026.

\bibitem{antlr-cpp}
ANTLR. \emph{The C++ Target}. Documentación del proyecto.
\url{https://github.com/antlr/antlr4/blob/dev/doc/cpp-target.md}.
Consulta: 31 de octubre de 2026.

\bibitem{flask}
Pallets Projects. \emph{Flask: Quickstart}.
\url{https://flask.palletsprojects.com/en/stable/quickstart/}.
Consulta: 31 de octubre de 2026.

\bibitem{subconsultas}
PostgreSQL Global Development Group. \emph{Subquery Expressions}.
Referencia documental sobre EXISTS, IN y NOT IN.
\url{https://www.postgresql.org/docs/current/functions-subquery.html}.
Consulta: 31 de octubre de 2026. No se utiliza PostgreSQL en el sistema.

\bibitem{repositorio}
Gonzalez, L. y Jaquez, J. \emph{Intérprete SQL}, rama
\texttt{Compiladores2-Proyecto3}, versión \texttt{037ac10}.
\url{https://github.com/G1LB3T0/Proyecto1_Compis_Generador_Analizador_Lexico/tree/Compiladores2-Proyecto3}.
Fuentes, ejemplos, pruebas y \path{docs/proyecto3-interprete-sql.md}.
\end{thebibliography}
\end{document}
